\subsection{Tables}
\label{subsec:impl-tables}

We distinguish between a table's \emph{local state}, maintained independently
by each node in the cluster and might differ among them, and its \emph{global
state}, derived from the combination of the local states over all nodes.

Locally, a table is in one of four phases: \emph{unknown}, \emph{prepared},
\emph{created}, or \emph{deleted}:
\begin{itemize}
    \item \emph{unknown}: the node is not aware of the table.
    \item \emph{prepared}: some throughput is allocated, but the table has not
        yet been created.
    \item \emph{created}: the table is fully operational on that node.
    \item \emph{deleted}: the table is no longer operational on that node.
\end{itemize}
The local state of a table can only progress in phase monotonically in the order:
\[
    \emph{unknown} \rightarrow \emph{prepared} \rightarrow \emph{created}
    \rightarrow \emph{deleted}.
\]

Globally, a table is in one of three phases: \emph{unknown}, \emph{created} or
\emph{deleted}.
A table cannot be globally \emph{prepared}, as this state represents an
intermediate phase during creation.
\begin{itemize}
    \item \emph{deleted}: at least one node is locally \emph{deleted}.
    \item \emph{created}: all responsible nodes for the table are locally
        \emph{created} and the table is not globally \emph{deleted}.
    \item \emph{unknown}: it is neither globally \emph{created} nor globally
        \emph{deleted}.
\end{itemize}
The global state of a table can also only progress in phase monotonically in the order:
\[
    \emph{unknown} \rightarrow \emph{created} \rightarrow \emph{deleted}.
\]

\subsubsection{Gossip}
\label{subsubsec:gossip}

The local state of a table is synchronized among nodes through the use of a
Gossip protocol.
This allows the local state of a table to converge towards the global state,
but should not change the global state of a table.

To achieve this, each node in the cluster initiates a \emph{gossip} exchange
with a randomly chosen peer every second.
If the peer is unavailable, the exchange fails and the node will try again in
the next second with a different peer.

During the exchange, each node tries to advertise all tables that are in the
locally \emph{created} or locally \emph{deleted} state to the peer.
Tables that are locally \emph{prepared} are not advertised, as this state is
only a local phase and not a global one.
A peer accepts the new state of a table only if it is in a more advanced phase
than its local state with respect to the order \emph{deleted} $\succ$
\emph{created} $\succ$ \emph{prepared} $\succ$ \emph{unknown}.
If the advertised state is in the same phase as the local state, then both
states must be the same; otherwise the system is inconsistent.

In the exceptional case that the peer is advertising a \emph{created} table
that is locally \emph{prepared}, the node will refuse the advertisement as this
might change the global state from \emph{unknown} to \emph{created}.

Merkle trees~\cite{MerkleTree} are used to efficiently detect divergent tables,
allowing the nodes to advertise only those.

\subsubsection{Creation}
\label{subsubsec:creation}

To create a table, the following steps are performed:
\begin{enumerate}
    \item The client connects to a node of its choice, which acts as a
        \emph{router}.
    \item The client sends the table's parameters $B$, $N$,
        $R$, and $W$ to the router.
    \item The router generates an \texttt{id} for the table and computes the
        effective throughput $S = B \cdot N$ to account for replication.
    \item\label{step:prepare} The router iterates over the nodes in the
        cluster, skipping unavailable ones, requesting a throughput allocation
        for the new table, allocating up to $B$ on each AZ\@.
        Each node can allocate as much throughput as it has available, up to the
        requested amount, after allocating, the node marks the
        table as locally \emph{prepared}.
    \item Once sufficient throughput is obtained (at least $S$), the router
        defines the list of responsible nodes and their respective throughputs
        $s_i$.
    \item\label{step:commit} The router asks the nodes that responded in
        step~\ref{step:prepare} to commit the table; on these nodes the table
        enters the \emph{created} state.
    \item The router returns the \texttt{id} of the new table to the client.
\end{enumerate}

If all steps succeed, the table is created successfully.
In particular if step~\ref{step:commit} succeeds, all the responsible nodes
will have marked the table as locally \emph{created}, and thus the table will
be considered globally \emph{created}.

If any node fails to respond in step~\ref{step:commit}, the router will
consider the table creation failed, and will return an error to the client.
It will also inform the nodes that responded in step~\ref{step:prepare} to
delete the table, allowing the throughput to be freed.

If the router fails between step~\ref{step:prepare} and step~\ref{step:commit},
some nodes will have marked the table as locally \emph{prepared}, wasting the
reserved throughput.
To avoid this, each node that participates in the creation of a table maintains
a timeout of 1 minute.
If the node does not receive a commit request within this time, it will mark
the table as locally \emph{deleted}, freeing the allocated throughput.

\subsubsection{Deletion}
\label{subsubsec:deletion}

To delete a table, it is sufficient for a node to mark it as locally
\emph{deleted}.
The table will then be considered globally \emph{deleted} and, thanks to the
Gossip protocol, the other nodes will synchronize and the table will be
considered deleted by all nodes.

Even if it is not strictly necessary, the node receiving the deletion request
acts as a router, informing all responsible nodes of the deletion, to allow
them to free the allocated throughput and storage space more quickly, and
prevent further read or write requests to the deleted table.
